\subsection*{Cross-model comparison on matched inputs}
\label{sec:crossmodel}

The two atlases were built from different cells, so any direct comparison of their statistics
confounds architecture with input distribution. To separate the two we trained Geneformer
dictionaries on the same 3,000 Tabula Sapiens cells used for the scGPT atlas, with identical
architecture and hyperparameters, and compared the models at matched relative depth
(\suptab{crossmodel}). Because the scGPT dictionaries were originally trained on all available
positions rather than a subsample, we also retrained them under the Geneformer protocol as a
control.

On matched inputs there is no uniform architectural advantage. Geneformer reconstructs better at
the input end (0.7731 against 0.7258 at relative depth 0), the two are indistinguishable at about a
quarter depth (0.7749 against 0.7729), and scGPT is ahead through the second half, reaching 0.8862
against 0.7241 at the output end. The difference between the models is a depth-dependent crossover
rather than a constant offset, and it is invisible unless the inputs are matched.

Input distribution accounts for a difference of the same order. Each Geneformer dictionary
evaluated on the distribution it was trained on reconstructs the CRISPRi-derived control cohort
4.7--7.8 percentage points better than the Tabula Sapiens dictionary reconstructs Tabula
Sapiens---0.8444 against 0.7731 at layer~0, for example---so a difference of this size between two
models is what differing inputs alone can produce. Both dictionaries sit close to the packing
limit for their dimensionality: mean absolute decoder coherence is 0.030--0.036 for Geneformer and
0.039--0.047 for scGPT, against Welch bounds of 0.0255 and 0.0383 respectively.

\subsection*{Robustness to training choices}
\label{sec:robustness}

\paragraph{Architecture.} Grid-training at layer~11 over expansion ratio $\{2\times, 4\times,
8\times\}$ and sparsity $k \in \{16, 32, 64\}$ gives variance explained rising monotonically with
both, from 0.706 to 0.884, with the reported configuration in the middle at 0.819
(\suptab{hyperparam}). Module count stays between 6 and 12 across the entire grid and mean decoder
coherence between 0.033 and 0.042, so the reported structure is not a property of one corner of the
hyperparameter space.

\paragraph{Random seed.} The atlas dictionaries were trained with seed 42. Retraining at layers 0
and 11 under six runs each---five differing only in initialisation and batch order, one with a
different training subsample---separates what is a property of the model from what is a property of
the run (\suptab{seed_stability}).

Everything the atlas reports in aggregate is reproducible. Variance explained is
$0.8454 \pm 0.0003$ at layer~0 and $0.8200 \pm 0.0003$ at layer~11, the annotation rate is
$0.935 \pm 0.009$ and $0.933 \pm 0.010$, module counts are $7.2 \pm 0.4$ and $8.3 \pm 1.0$, and
dead-feature counts are $2 \pm 1$ and $32 \pm 4$.

Individual dictionary atoms are not. The mean best decoder cosine between runs is 0.282 at layer~0
and 0.294 at layer~11, and fewer than 1\% of atoms have a counterpart above 0.9 in another run. The
reference for these numbers is what a random direction achieves: the best match a random unit
direction finds among 4,608 directions in 1,152 dimensions is $0.113 \pm 0.008$, and no random
direction reaches 0.9. Cross-seed agreement is therefore well above chance and far below
reproducibility. Independent runs learn different bases that explain the same variance and carry the
same annotation. This bears on how the released atlases should be read: their aggregate structure
is a property of the model, while the identity of any individual catalogued feature is partly a
property of the training run---consistent with the gene-level analysis above, in which the recurring
unit across layers is the program rather than the atom.

\subsection*{Features without ontology annotation}
\label{sec:unannotated}

Between 41\% and 55\% of features carry no significant ontology annotation, and it matters whether
these are structure the databases do not cover or noise. Clustering unannotated features by
top-20 gene-set similarity places only 2--3\% into standalone gene-set clusters, among them
ribosomal protein programs and mitochondrial complex assembly (\suptab{novel}).

Two tests separate them from annotated features (\suptab{unannotated_new}). Cell-type specificity
does not: at layers $\{0, 5, 11, 17\}$, 89.4--95.0\% of unannotated features are enriched for at
least one cell type against 94.5--95.8\% of annotated ones, and a matched random-feature null sits
at 91.9--94.7\%, so essentially every feature at these layers is cell-type selective and the test
has no discriminating power. Perturbation responsiveness does: at layer~11, 2.32\% of unannotated
features appear in a perturbation response list against 0.15\% of annotated features, a fifteenfold
difference, with unannotated features more responsive at layers 0 and 5 and indistinguishable at layer 17.

A cautious reading is that some unannotated features carry perturbation-response signal the
ontology databases do not describe. Being non-random is not the same as being biologically
meaningful, however: a feature that is cell-type selective may equally be tracking donor or assay
structure, which is why the batch analysis above is reported alongside.
